\documentclass[11pt]{article}
% Shared preamble of the RMC kernel write-ups (% Shared preamble of the RMC kernel write-ups (% Shared preamble of the RMC kernel write-ups (\input{preamble} after \documentclass).
\usepackage[margin=1in]{geometry}
\usepackage{amsmath,amssymb,bm}
\usepackage{graphicx}
\usepackage{booktabs}
\usepackage{hyperref}

\newcommand{\dd}{\,\mathrm{d}}
\newcommand{\Ein}{E}
\newcommand{\Eout}{E'}
\newcommand{\Ecm}{E_{c}}
\newcommand{\mucm}{\mu_{c}}
\newcommand{\mulab}{\mu_0}

% Common closing section: how the verification figures are produced
\newcommand{\verificationmethod}{%
Each figure is produced by \texttt{verification/verify\_laws.py}. For an incident
energy $\Ein$ along $+z$, $10^6$ collisions are sampled with MC/DC's own reaction
sampler (\texttt{sample\_elastic\_scattering}, \texttt{sample\_inelastic\_scattering}
or \texttt{sample\_fission}), and every emitted neutron is histogrammed in
$48 \times 24$ bins of $(\Eout, \mulab)$, with $\mulab$ the lab cosine to $+z$. The
inverted kernel used by RMC is integrated over the same bins (Gauss--Legendre panels;
two-body lines are split where $\mulab^*(\Eout)$ crosses a cosine edge), multiplied by
the number of collisions, and compared with a $\chi^2$ test over bins with at least 5
expected counts (the rest pooled into one bin). The tables list $\chi^2/\mathrm{dof}$,
its $p$-value, the emitted neutrons per collision (sampled vs.\ the integral of the
inverted kernel), and the largest relative change of a bin integral when the
quadrature panels are doubled. Left panels: outgoing-energy marginal; middle:
lab-cosine marginal; right: per-bin $z = (\text{sampled} - \text{inverted}) /
\sqrt{\text{inverted}}$.}
 after \documentclass).
\usepackage[margin=1in]{geometry}
\usepackage{amsmath,amssymb,bm}
\usepackage{graphicx}
\usepackage{booktabs}
\usepackage{hyperref}

\newcommand{\dd}{\,\mathrm{d}}
\newcommand{\Ein}{E}
\newcommand{\Eout}{E'}
\newcommand{\Ecm}{E_{c}}
\newcommand{\mucm}{\mu_{c}}
\newcommand{\mulab}{\mu_0}

% Common closing section: how the verification figures are produced
\newcommand{\verificationmethod}{%
Each figure is produced by \texttt{verification/verify\_laws.py}. For an incident
energy $\Ein$ along $+z$, $10^6$ collisions are sampled with MC/DC's own reaction
sampler (\texttt{sample\_elastic\_scattering}, \texttt{sample\_inelastic\_scattering}
or \texttt{sample\_fission}), and every emitted neutron is histogrammed in
$48 \times 24$ bins of $(\Eout, \mulab)$, with $\mulab$ the lab cosine to $+z$. The
inverted kernel used by RMC is integrated over the same bins (Gauss--Legendre panels;
two-body lines are split where $\mulab^*(\Eout)$ crosses a cosine edge), multiplied by
the number of collisions, and compared with a $\chi^2$ test over bins with at least 5
expected counts (the rest pooled into one bin). The tables list $\chi^2/\mathrm{dof}$,
its $p$-value, the emitted neutrons per collision (sampled vs.\ the integral of the
inverted kernel), and the largest relative change of a bin integral when the
quadrature panels are doubled. Left panels: outgoing-energy marginal; middle:
lab-cosine marginal; right: per-bin $z = (\text{sampled} - \text{inverted}) /
\sqrt{\text{inverted}}$.}
 after \documentclass).
\usepackage[margin=1in]{geometry}
\usepackage{amsmath,amssymb,bm}
\usepackage{graphicx}
\usepackage{booktabs}
\usepackage{hyperref}

\newcommand{\dd}{\,\mathrm{d}}
\newcommand{\Ein}{E}
\newcommand{\Eout}{E'}
\newcommand{\Ecm}{E_{c}}
\newcommand{\mucm}{\mu_{c}}
\newcommand{\mulab}{\mu_0}

% Common closing section: how the verification figures are produced
\newcommand{\verificationmethod}{%
Each figure is produced by \texttt{verification/verify\_laws.py}. For an incident
energy $\Ein$ along $+z$, $10^6$ collisions are sampled with MC/DC's own reaction
sampler (\texttt{sample\_elastic\_scattering}, \texttt{sample\_inelastic\_scattering}
or \texttt{sample\_fission}), and every emitted neutron is histogrammed in
$48 \times 24$ bins of $(\Eout, \mulab)$, with $\mulab$ the lab cosine to $+z$. The
inverted kernel used by RMC is integrated over the same bins (Gauss--Legendre panels;
two-body lines are split where $\mulab^*(\Eout)$ crosses a cosine edge), multiplied by
the number of collisions, and compared with a $\chi^2$ test over bins with at least 5
expected counts (the rest pooled into one bin). The tables list $\chi^2/\mathrm{dof}$,
its $p$-value, the emitted neutrons per collision (sampled vs.\ the integral of the
inverted kernel), and the largest relative change of a bin integral when the
quadrature panels are doubled. Left panels: outgoing-energy marginal; middle:
lab-cosine marginal; right: per-bin $z = (\text{sampled} - \text{inverted}) /
\sqrt{\text{inverted}}$.}

\usepackage{amsthm}

\newtheorem{proposition}{Proposition}
\newcommand{\tpsi}{\tilde\psi}
\newcommand{\Sbar}{\bar S}
\newcommand{\Sigt}{\Sigma_t}
\newcommand{\E}{\mathbb{E}}
\newcommand{\Var}{\operatorname{Var}}

\title{Sampling the scattering correction: pointwise and integrated estimators}
\date{}

\begin{document}
\maketitle

\noindent Notation as in \texttt{rmc\_math.tex}: an outgoing phase-space point
$x=(z,E,\mu)$ lies in a trial-space bin $\beta=(k,g,j)$; $E'$ is the incident energy,
in an incident bin $g'$. Section~6 gives the extension to the linear trial spaces.

\section{The quantity to sample}

At the fixed point of the collision-only iterations, each correction pass samples the
full residual. Its scattering part (fission is identical) is
\begin{equation}
  r_s(x) = T(x) - \Sbar_\beta(x),\qquad
  T(x) = \int \dd E'\, t(E',x),
  \label{eq:rs}
\end{equation}
where
\begin{equation}
  t(E',x) = \sum_{j'} \tpsi_{kg'j'}(E')\, U_{j'}(E'\to E,\mu),\qquad
  U_{j'} = \sum_r N_r\sigma_r(E')\,\bar\tau_{r,j'}(E'\to E,\mu),
\end{equation}
is the in-scatter density per unit incident energy (the incident polar cosine is
already integrated over each bin $j'$ in closed form), and $\Sbar_\beta$ is its
projection onto the trial space. By construction $\int_\beta r_s\dd x = 0$: $r_s$
carries only the in-bin shape of the scattering source. Both samplers below estimate
the transport response to $r_s$ by banking particles at sampled points with signed
weights $w = r/q$, the standard importance-sampling identity for a signed source [1, 2].

\section{The two estimators}

\paragraph{Pointwise.} Sample the joint point $y=(E',x)$ from a density $q(y)$ and score
\begin{equation}
  w_{\rm pw}(y) = \frac{t(E',x) - \bar s(E',x)}{q(y)},
  \qquad \int \bar s(E',x)\dd E' = \Sbar_\beta(x),
  \label{eq:pointwise}
\end{equation}
with a subtracted density $\bar s$ per unit incident energy. The implementation takes
$\bar s$ flat in $E'$ over each incident bin:
$\bar s_{g'} = \Sbar^{(g')}_\beta/\Delta E_{g'}$, the contribution of bin $g'$ to
$\Sbar_\beta$ spread uniformly.

\paragraph{Integrated.} Sample only $x$ from $q_x(x)$ and do the $E'$ integral
deterministically:
\begin{equation}
  w_{\rm int}(x) = \frac{T(x) - \Sbar_\beta(x)}{q_x(x)} = \frac{r_s(x)}{q_x(x)} .
  \label{eq:integrated}
\end{equation}

\begin{proposition}[Unbiasedness]
For any response function $\phi(x)$ (the transport of a particle born at $x$, in
expectation), and proposals positive wherever the scored numerator is nonzero,
\[
  \E_q[w_{\rm pw}\,\phi(x)] = \iint (t - \bar s)\,\phi\dd E'\dd x
  = \int r_s\,\phi\dd x = \E_{q_x}[w_{\rm int}\,\phi(x)] .
\]
\end{proposition}
\noindent Both estimators are unbiased for the exact $T$. The pointwise one needs no
quadrature, so it remains unbiased as implemented. The integrated one is unbiased only
up to the accuracy of the computed $\widehat T$ (Section~4).

\section{Second moments}

\paragraph{Lower bounds.} By the Cauchy--Schwarz inequality, over all proposals,
\begin{equation}
  \E[w_{\rm pw}^2] = \iint\frac{(t-\bar s)^2}{q}
  \;\ge\; \Big(\iint|t-\bar s|\dd E'\dd x\Big)^2 \equiv L_{\rm pw}^2,
  \qquad
  \E[w_{\rm int}^2] \;\ge\; \Big(\int|r_s|\dd x\Big)^2 = \|r_s\|_1^2,
\end{equation}
with equality for $q\propto|t-\bar s|$ and $q_x\propto|r_s|$. By the triangle
inequality, $|r_s(x)| = |\int(t-\bar s)\dd E'| \le \int|t-\bar s|\dd E'$, so
$\|r_s\|_1 \le L_{\rm pw}$.

\begin{proposition}[Rao--Blackwell]
Let the pointwise proposal factor as $q(y) = q_x(x)\,q(E'\mid x)$, with the same $q_x$
as the integrated one. Then $\E[w_{\rm pw}\mid x] = w_{\rm int}(x)$ and
\[
  \Var(w_{\rm pw}) = \Var(w_{\rm int}) + \E\big[\Var(w_{\rm pw}\mid x)\big]
  \;\ge\; \Var(w_{\rm int}).
\]
\end{proposition}
\begin{proof}
$\E[w_{\rm pw}\mid x] = \int \frac{t-\bar s}{q_x\,q(E'|x)}\,q(E'|x)\dd E'
= (T-\Sbar_\beta)/q_x$; then apply the law of total variance.
\end{proof}
\noindent This is conditional Monte Carlo, or Rao--Blackwellization [3--5]: integrating
out a variable analytically never increases variance. The excess,
$\E[\Var(w_{\rm pw}\mid x)]$, is the variance of the incident-energy integrand at fixed
$x$.

\paragraph{Why the excess is large.} The incident-energy integrand $t(E',x)$ is not flat.
For elastic scattering (a two-body law), an outgoing $E$ is reached only from
$E'\in[E, E/\alpha]$, $\alpha = ((A-1)/(A+1))^2$. On the diagonal blocks $g'\approx g$,
$t$ is a triangle in $(E',E)$ that vanishes at the kinematic edge, while the flat
$\bar s_{g'}$ does not. Then $|t-\bar s|$ is of the size of $t$ itself, so
$L_{\rm pw}$ scales with the full in-scatter source, while $\|r_s\|_1$ is only its
in-bin shape. The ratio of the attainable second moments,
\begin{equation}
  R = \frac{L_{\rm pw}^2}{\|r_s\|_1^2} \;\ge\; 1,
\end{equation}
grows as the trial space improves, because $\|r_s\|_1\to 0$ while $L_{\rm pw}$ stays
of the order of the scattering source. On the $A=1$ analytic problem with $G=100$
energy bins, the pointwise sampler's measured $\E[w^2]$ was about $485\times$ the ideal
$\|r_s\|_1^2$ (\texttt{rmc\_math.tex}). Improving the trial space shrinks the
integrated estimator's variance and leaves the pointwise one's nearly unchanged.

\section{Bias of the computed integral}

The implementation evaluates $T$ by adaptive Gauss--Kronrod quadrature in $E'$, giving
$\widehat T$. The banked source is then $\widehat T - \Sbar_\beta$, whose expectation
differs from $r_s$ by the deterministic error $\widehat T - T$. Two conditions make that
error controllable:
\begin{itemize}
  \item \emph{Breakpoints at every non-smooth point of the integrand.} These are the
  trial-space edges (where $\tpsi$ jumps), the incident energies where a reaction's
  emission support crosses $E$, the reaction thresholds, and the cross-section grid
  points inside the kinematic window $W(E)=\{E' : E\in\operatorname{supp}
  f(E'\to\cdot)\}$. Between grid points $\sigma_r(E')$ is linear, so the integrand is
  smooth on each piece and Gauss rules converge rapidly. Without the cross-section
  points, a narrow resonance that falls between all quadrature nodes is missed,
  leaving a deterministic error in resonance bins. Kinematic pruning keeps the number
  of points small: for elastic scattering $W(E)=[E, E/\alpha]$, a few percent of the
  grid for heavy nuclides.
  \item \emph{An absolute error scale.} $r_s$ is a small difference of two large terms,
  so a relative tolerance on $\widehat T$ alone does not bound the error relative to
  $\|r_s\|_1$. The tolerance should be set against $\|r_s\|_1$ (or the expected
  statistical error of the pass), not against $|T|$.
\end{itemize}
The kinematic breakpoints depend only on the reaction and on $E$, not on $\tpsi$ or on
the sample, so they can be tabulated once per outgoing bin.

\section{Efficiency}

The figure of merit of a correction pass is $\mathrm{FOM} = 1/(\sigma^2\tau)$, with
$\tau$ the cost per history. Each history costs its source sampling $\tau_{\rm src}$
plus its transport $\tau_{\rm tr}$, and the transport part is the same for both
samplers. Their efficiency ratio is
\begin{equation}
  \frac{\mathrm{FOM}_{\rm int}}{\mathrm{FOM}_{\rm pw}}
  = \frac{\E[w_{\rm pw}^2]}{\E[w_{\rm int}^2]}\cdot
    \frac{\tau^{\rm pw}_{\rm src}+\tau_{\rm tr}}{\tau^{\rm int}_{\rm src}+\tau_{\rm tr}}
  \;\approx\; R\cdot\frac{\tau^{\rm pw}_{\rm src}+\tau_{\rm tr}}
  {n_{\rm eval}\,\tau_{\rm ker}+\tau_{\rm tr}},
  \label{eq:fom}
\end{equation}
(both weights have the same mean, so the ratio of variances approaches that of second
moments when, as here, the spread of the signed weights dominates their mean), where
$n_{\rm eval}$ is the number of kernel evaluations in one adaptive $E'$ integral and
$\tau_{\rm ker}$ the cost of one evaluation, summed over the reactions. Two limits:
\begin{itemize}
  \item \emph{Transport-dominated} ($\tau_{\rm tr}\gg n_{\rm eval}\tau_{\rm ker}$, e.g.\
  thick, highly scattering media where each history makes many collisions): the ratio
  tends to $R$, and the integrated sampler wins by the full variance ratio.
  \item \emph{Source-dominated} ($n_{\rm eval}\tau_{\rm ker}\gg\tau_{\rm tr}$, e.g.\
  many reactions, wide kinematic windows, short histories): the ratio tends to
  $R/n_{\rm eval}$. The integrated sampler wins only if one integral costs fewer kernel
  evaluations than the variance ratio $R$.
\end{itemize}
Since $R$ grows as the trial space improves while $n_{\rm eval}$ does not, the balance
shifts toward the integrated estimator as the method matures. Tabulating the
breakpoints (Section~4) reduces $n_{\rm eval}\tau_{\rm ker}$ directly.

\section{Linear trial spaces}

With the energy, angular and spatial trial functions of \texttt{linear\_energy\_basis.tex},
\texttt{linear\_discontinuous\_angle.tex} and \texttt{continuous\_linear\_space.tex},
the in-scatter density at $x=(z,E,\mu)$ in cell $k$ is
\begin{equation}
  T(x) = \sum_{g'j'}\sum_{a'b'} \tpsi_{g'j'a'b'}(z)\int_{g'}\dd E'\,
  P_{a'}\big(x_{g'}(E')\big) \sum_r N_r\sigma_r(E')\,
  \bar\tau^{\,b'}_{r,j'}(E'\to E,\mu),
\end{equation}
\begin{equation}
  \bar\tau^{\,b'}_{r,j'} = \int\dd\mu_0\, f_{r,L}(E',E,\mu_0)\,
  \Pi^{b'}_{j'}(\mu,\mu_0),
\end{equation}
where $\tpsi_{\cdot}(z)$ interpolates the nodal coefficients linearly for the
continuous spatial basis, and $\Pi^{b'}$ is the closed-form angular moment of the
incident bin (the azimuthal kernel is symmetric in the incident and outgoing cosines).
The projection
\[
  \Sbar_\beta(x) = \sum_{a,b}\frac{(2a+1)(2b+1)}{\Delta E_g\Delta\mu_j}\,
  P_a\big(x_g(E)\big)P_b\big(y_j(\mu)\big)\sum M\,\tpsi
\]
is linear in $E$ and $\mu$, and linear in $z$ for the continuous spatial basis.
Both estimators carry over: the integrated one evaluates $T(x)$ and $\Sbar_\beta(x)$ at
the sampled $x$; the pointwise one needs a subtracted density with
$\int\bar s\dd E' = \Sbar_\beta(x)$ (for instance flat in $E'$ over each $g'$, as
above, times the polynomial factors in $E$ and $\mu$). The variance argument of
Section~3 is unchanged: a better trial space reduces $\|r_s\|_1$ and therefore the
integrated estimator's variance, but not $L_{\rm pw}$.

\section{A family between the two: randomized quadrature in $E'$}

Draw $x\sim q_x$, and then $m$ incident energies $E'_1,\dots,E'_m$ by stratified
sampling over the pieces between breakpoints: one point uniform in each of $m$
equal-probability strata of a density $q(E'\mid x)$. Score
\begin{equation}
  w_m(x) = \frac{1}{q_x(x)}\,\frac{1}{m}\sum_{i=1}^m
  \frac{t(E'_i,x) - \bar s(E'_i,x)}{q(E'_i\mid x)} .
\end{equation}
\begin{proposition}
$w_m$ is unbiased for every $m\ge 1$, with $\E[w_m\mid x] = w_{\rm int}(x)$. Its excess
variance $\E[\Var(w_m\mid x)]$ is at most that of $m$ independent draws, i.e.\ the
pointwise excess divided by $m$. For an integrand with a bounded derivative on each
stratum it falls as $O(m^{-3})$ [5].
\end{proposition}
\noindent For $m=1$ with a single stratum this is the pointwise sampler, and as
$m\to\infty$ it tends to the integrated one. For finite $m$ it is exactly unbiased: it
needs no quadrature tolerance and has no deterministic error at missed resonances.
Its cost is $m$ kernel evaluations per particle, so $m$ can be chosen to balance
\eqref{eq:fom}, typically a few strata per piece between breakpoints. It is a natural
candidate to replace both estimators.

\paragraph{Support of the proposals.} Unbiasedness needs $q>0$ wherever the scored
numerator can be nonzero. Pilot masses estimated at a few points can vanish where
$r_s$ does not (a signed integrand can be zero at every pilot node). A defensive
mixture $q_x = (1-\epsilon)\,q_{\rm pilot} + \epsilon\, q_{\rm vol}$, with $q_{\rm vol}$
proportional to bin volume, guarantees support for any $\epsilon>0$, and costs at most
a factor $1/(1-\epsilon)$ in the second moment where the pilot is good [6].

\section{Summary}

\begin{center}
\begin{tabular}{lll}
\toprule
 & pointwise & integrated \\
\midrule
Bias & none & quadrature error of $\widehat T$ \\
Best second moment & $L_{\rm pw}^2$ (in-scatter scale) & $\|r_s\|_1^2$ (residual scale) \\
Improves with the trial space & no & yes \\
Cost per particle & one kernel set & one adaptive $E'$ integral \\
\bottomrule
\end{tabular}
\end{center}

\noindent With the measured $R\approx 485$ on the $A=1$ problem, the integrated
estimator is more efficient whenever one $E'$ integral costs fewer than about $485$
kernel evaluations, or whenever transport dominates the cost per history. The
stratified estimator of Section~7 keeps most of that variance reduction without any
quadrature bias.

\begin{thebibliography}{99}
\bibitem{spanier} J.~Spanier and E.~M.~Gelbard, \emph{Monte Carlo Principles and
  Neutron Transport Problems}, Addison-Wesley (1969).
\bibitem{halton} J.~H.~Halton, ``Sequential Monte Carlo,'' \emph{Proc.\ Cambridge
  Philos.\ Soc.} \textbf{58}, 57 (1962).
\bibitem{hammersley} J.~M.~Hammersley and D.~C.~Handscomb, \emph{Monte Carlo
  Methods}, Methuen (1964), Ch.~5 (conditional Monte Carlo).
\bibitem{blackwell} D.~Blackwell, ``Conditional expectation and unbiased sequential
  estimation,'' \emph{Ann.\ Math.\ Statist.} \textbf{18}, 105 (1947); C.~R.~Rao,
  ``Information and the accuracy attainable in the estimation of statistical
  parameters,'' \emph{Bull.\ Calcutta Math.\ Soc.} \textbf{37}, 81 (1945).
\bibitem{owen} A.~B.~Owen, \emph{Monte Carlo Theory, Methods and Examples} (2013),
  chapters on conditioning, stratification and importance sampling.
\bibitem{hesterberg} T.~Hesterberg, ``Weighted average importance sampling and
  defensive mixture distributions,'' \emph{Technometrics} \textbf{37}, 185 (1995).
\bibitem{vermaak} J.~I.~C.~Vermaak and J.~E.~Morel, ``Transport Error Estimation Using
  Residual Monte Carlo,'' \emph{J.\ Comput.\ Phys.} (2022).
\bibitem{larsen} M.~Larsen et al., ``A Residual Monte Carlo Algorithm for Continuous
  Energy Neutron Transport with Elastic Scattering,'' \emph{Nucl.\ Sci.\ Eng.}
  \textbf{200}(3) (2025).
\end{thebibliography}

\end{document}
